\documentclass[../../../uem_phy_oka_q.tex]{subfiles}

\begin{document}
    図\ref{fig:uem_phy_oka_2025_uni_03_q_01}のように，%
    容器AとBがバルブのついた細い管で接続されている．容器Aは，断面積$S$のシリンダーと%
    なめらかに動くピストンからなる．容器Bは断面積$S$で長さ$\myfraca{L}{4}$である．%

    最初にバルブは閉じられている．容器Aには単原子分子の理想気体$n$モルが閉じ込められており，%
    ピストンはシリンダーの右端から長さ$L$の位置にストッパーが固定され，%
    内部の圧力は大気圧$P_0$と等しい．また，容器Bの内部は真空である．%
    これを状態\mynumberb{1}とする．気体定数を$R$として以下の問いに答えよ．%
    ただし，容器A，Bとバルブおよび管は断熱材で作られていて，%
    ストッパーの体積は無視できるものとする．

    \begin{enumerate}
        \item%
            状態\mynumberb{1}における気体の温度$T_1$と内部エネルギーを求めよ．%
            $\{P_0,\,S,\,L,\,n,\,R\}$
    \end{enumerate}

    状態\mynumberb{1}から，バルブを静かに開き，充分時間が経過すると全体の状態が一様になった．%
    このとき，気体の温度は$T_2$であった．この状態を状態\mynumberb{2}とする．

    \begin{enumerate}[resume]
        \item%
            状態\mynumberb{1}における気体の温度$T_1$と$T_2$の関係を，%
            解答欄の選択肢$\{<,\,=,\,>\}$の中から選んで丸で囲み，その理由を答えよ．
        \item%
            状態\mynumberb{2}で，ピストンを手に持ちストッパーを外したとき，%
            ピストンの位置を一定に保つには力を加える必要がある．%
            その力の大きさを求めよ．$\{P_0,\,S,\,L\}$
    \end{enumerate}

    あらためて，状態\mynumberb{1}から，ストッパーを外してからバルブを静かに開き，%
    充分時間が経過すると全体の状態が一様になった．このとき，ピストンは右側に動き，%
    その位置はシリンダーの右端から$L_3$となった．また，気体の温度は$T_3$であった．%
    この状態を状態\mynumberb{3}とする．%

    \begin{enumerate}[resume]
        \item%
            状態\mynumberb{1}から状態\mynumberb{3}へ変わる間に%
            気体が外部にした仕事を求めよ．$\{P_0,\,S,\,L,\,L_3\}$
        \item%
            状態\mynumberb{3}におけるピストンの位置$L_3$を$L$を用いて表せ．
        \item%
            状態\mynumberb{3}における気体の温度$T_3$を$T_1$を用いて表せ．
    \end{enumerate}
    \newpage
    \noindent
    \begin{minipage}{\linewidth}
        \centering
            \includegraphics{fig_uem_phy_oka_2025_uni_03/fig_uem_phy_oka_2025_uni_03_q_01.pdf}
            \captionof{figure}{} \label{fig:uem_phy_oka_2025_uni_03_q_01}
    \end{minipage}
\end{document}
